\chapter{Konstrukcja liczb rzeczywistych}
\label{ma1:ch:real-construction}
\textbf{Cel rozdziału.} Skonstruujemy ciało liczb rzeczywistych przez przekroje Dedekinda, a następnie jako zbiór klas ciągów Cauchy'ego liczb wymiernych. Pierwsza konstrukcja uzupełnia luki w porządku $\Q$, druga uzupełnia luki w zbieżności. Warunek Cauchy'ego sformułujemy najpierw za pomocą dodatnich liczb wymiernych, nie zakładając istnienia $\R$.

\section{Przekroje Dedekinda}
\label{ma1:sec:dedekind}
W rozdziale \ref{ma1:ch:number-systems} wykazaliśmy, że zbiór nieujemnych liczb wymiernych o kwadracie mniejszym od $2$ nie ma kresu górnego w $\Q$. Zamiast szukać brakującego kresu w dotychczasowym zbiorze, opiszemy nową liczbę przez wszystkie liczby wymierne leżące poniżej niej.

\begin{definicja}[Przekrój Dedekinda]
\label{ma1:def:cut}
Przekrojem Dedekinda nazywamy zbiór $A\subseteq\Q$ taki, że $A\ne\varnothing$, $A\ne\Q$, dla wszystkich $a\in A$ i $q\in\Q$ takich, że $q<a$, zachodzi $q\in A$, a dla każdego $a\in A$ istnieje $a'\in A$ takie, że $a<a'$. Zbiór wszystkich przekrojów oznaczamy przez $\R_D$. Jeśli $a\in A$ i $u\in\Q\setminus A$, to $a<u$: w przypadku $u=a$ lub $u<a$ mielibyśmy $u\in A$. W szczególności dopełnienie przekroju jest domknięte w górę.
\end{definicja}
Dla $r\in\Q$ piszemy $D_r=\{q\in\Q:q<r\}$. Jest to przekrój: zawiera $r-1$, nie zawiera $r$, jest domknięty w dół, a dla $q<r$ liczba $(q+r)/2$ należy do $D_r$ i jest większa od $q$.
\begin{twierdzenie}[Porządek i zupełność]
\label{ma1:thm:cut-complete}
Relacja $A\le B$ określona przez $A\subseteq B$ jest porządkiem liniowym na $\R_D$. Każda niepusta rodzina $\mathcal A\subseteq\R_D$ ograniczona z góry ma kres górny równy $\bigcup_{A\in\mathcal A}A$. Odwzorowanie $r\mapsto D_r$ zachowuje i odzwierciedla porządek. Dla wszystkich $A,B\in\R_D$ takich, że $A<B$ (tj. $A\subsetneq B$), istnieje $r\in\Q$ takie, że $A<D_r<B$.
\end{twierdzenie}
\begin{proof}
Zwrotność, antysymetria i przechodniość wynikają z własności inkluzji. Jeżeli $A\not\subseteq B$, wybierzmy $a\in A\setminus B$. Każde $b\in B$ spełnia $b<a$, gdyż w przeciwnym razie $a\in B$; zatem $B\subseteq A$. Jeśli $U$ jest ograniczeniem górnym rodziny $\mathcal A$, jej suma zbiorów $V$ jest niepusta i zawiera się w $U\ne\Q$. Domknięcie w dół i brak elementu największego wynikają z tych własności każdego składnika. Każdy składnik zawiera się w $V$, a każde ograniczenie górne zawiera $V$. Stąd $V=\sup\mathcal A$.
Dla $r<s$ mamy $D_r\subsetneq D_s$, ponieważ $r\in D_s\setminus D_r$. Liniowość porządku $\Q$ daje odzwierciedlanie porządku. Jeśli $A<B$, wybieramy $a\in B\setminus A$ oraz $b\in B$ z $b>a$. Dla $a<r<b$ wszystkie elementy $A$ są mniejsze od $a$, natomiast $a\in D_r\setminus A$ oraz $b\in B\setminus D_r$. Mamy więc $A<D_r<B$.
\end{proof}
\begin{lemat}[Przybliżenia wymierne przekroju]
\label{ma1:lem:cut-bracket}
Dla $A\in\R_D$ i $\varepsilon\in\Q$, $\varepsilon>0$, istnieją $a\in A$, $b\in\Q\setminus A$ takie, że $0<b-a<\varepsilon$.
\end{lemat}
\begin{proof}
Wybierzmy $u\in A$, $v\notin A$; wówczas $u<v$. Połóżmy $h=\varepsilon/2$. Własność Archimedesa w $\Q$ zapewnia istnienie $n\in\N$ z $nh>v-u$. Zbiór $\{k\in\N:u+kh\notin A\}$ jest więc niepusty. Jego najmniejszy element $k$ jest dodatni. Wystarczy przyjąć $a=u+(k-1)h$ oraz $b=u+kh$.
\end{proof}

\section{Dodawanie i element przeciwny}
\label{ma1:sec:cut-addition}
\begin{definicja}
Dla $A,B\in\R_D$ definiujemy
\[
 A+B=\{a+b:a\in A,\ b\in B\},\qquad
 -A=\{q\in\Q:\exists u\notin A,\ q<-u\}.
\]
Przyjmujemy $0_D=D_0$, $1_D=D_1$.
\end{definicja}
\begin{twierdzenie}
\label{ma1:thm:cut-group}
Zbiór $\R_D$ z dodawaniem jest grupą abelową z zerem $0_D$; elementem przeciwnym do $A$ jest $-A$. Porządek jest zgodny z dodawaniem, a $D_r+D_s=D_{r+s}$ dla wszystkich $r,s\in\Q$.
\end{twierdzenie}
\begin{proof}
Suma jest niepusta. Jeśli $u\notin A$, $v\notin B$, to każdy $a+b<u+v$, więc suma jest właściwa. Dla $q<a+b$ mamy $q-b<a$, zatem $q=(q-b)+b$ należy do sumy. Zwiększenie $a$ wewnątrz $A$ pokazuje brak elementu największego. Przekrój $-A$ jest niepusty (dla $u\notin A$ zawiera $-u-1$), właściwy (jeśli $a\in A$, to każde $q\in -A$ spełnia $q<-a$), domknięty w dół i bez elementu największego dzięki gęstości $\Q$.
Łączność i przemienność sum wynikają z tych własności w $\Q$. Każde $a+z$ dla $a\in A$, $z<0$ należy do $A$. Odwrotnie, dla $q\in A$ wybieramy $a>q$ w $A$, więc $q=a+(q-a)\in A+0_D$. Zatem $A+0_D=A$.
Jeżeli $a\in A$, $q\in -A$, to $q<-u$ dla pewnego $u\notin A$, a $a<u$. Stąd $a+q<0$. Aby dowieść odwrotnej inkluzji, niech $z<0$. Z lematu \ref{ma1:lem:cut-bracket} wybieramy $a\in A$, $u\notin A$ z $u-a<-z$. Wówczas $z-a<-u$, więc $z=a+(z-a)\in A+(-A)$. Otrzymujemy $A+(-A)=0_D$.
Inkluzja $A\subseteq B$ daje $A+C\subseteq B+C$; dodanie $-C$ dowodzi implikacji odwrotnej. Wreszcie $D_r+D_s\subseteq D_{r+s}$. Dla $q<r+s$ wybieramy $h=(r+s-q)/2>0$; wtedy $q=(r-h)+(s-h)$ należy do sumy. To kończy dowód.
\end{proof}

\section{Mnożenie}
\label{ma1:sec:cut-product}
\begin{definicja}[Iloczyn nieujemnych przekrojów]
Dla $A,B\in\R_D$ takich, że $A,B\ge0_D$, kładziemy
\[
 A\odot B=\{q\in\Q:q<0\ \text{lub istnieją }a\in A,\ b\in B,\ a,b>0,\ q<ab\}.
\]
\end{definicja}
\begin{lemat}
\label{ma1:lem:positive-products}
Iloczyn $\odot$ jest dobrze określoną operacją na nieujemnych przekrojach; jest przemienny, łączny, ma jedynkę $1_D$, zero pochłaniające $0_D$ i jest rozdzielny względem dodawania. Iloczyn dwóch przekrojów dodatnich jest dodatni.
\end{lemat}
\begin{proof}
Zauważmy, że suma dwóch nieujemnych przekrojów jest nieujemna: zgodność porządku z dodawaniem daje $0_D\le A\le A+B$. Zbiór iloczynu zawiera wszystkie liczby ujemne. Jeśli $u\notin A$, $v\notin B$ są dodatnie (możemy powiększyć dowolne elementy dopełnień), to $uv$ nie należy do iloczynu. Domknięcie w dół wynika z definicji. Dla $q<0$ wybieramy większą liczbę ujemną, a dla $q\ge0$ wybieramy liczbę między $q$ a świadkiem $ab$. To dowodzi, że iloczyn jest przekrojem nieujemnym. Gdy jeden czynnik to $0_D$, brak dodatnich świadków daje $0_D$. Gdy oba są dodatnie, zawierają dodatnie liczby $a,b$ (jeśli zawierają zero, brak maksimum daje liczbę dodatnią); wtedy $ab/2>0$ należy do iloczynu.
Przemienność jest bezpośrednia. Dla dodatnich $A,B,C$ i $q\ge0$ przynależność do $(A\odot B)\odot C$ jest równoważna istnieniu dodatnich $a\in A,b\in B,c\in C$ z $q<abc$. W jedną stronę używamy $q<uc$ i $0<u<ab$. W drugą wybieramy wymierne $u$ pomiędzy $q/c$ i $ab$, dodatnie. Analogiczna charakterystyka dla $A\odot(B\odot C)$ daje łączność; przypadek czynnika zerowego jest rozstrzygnięty. Dla $q\ge0$ z $q<ab$, $0<b<1$, wynika $q<a$, czyli $q\in A$. Jeżeli natomiast $q\in A$, $q\ge0$, wybieramy $a\in A$ większe od $q$ i $b\in\Q$ z $q/a<b<1$. Stąd $A\odot1_D=A$.
Dla dodatnich $A,B,C$ pokażemy rozdzielność. Jeśli $q\ge0$ należy do $A\odot(B+C)$, istnieją $a>0$ w $A$ oraz $s=b+c>0$ z $b\in B,c\in C$ i $q<as$. W każdym z dodatnich przekrojów $B,C$ istnieje dodatni element. Zastępując $b,c$ przez maksima z odpowiednimi dodatnimi elementami, otrzymujemy dodatnie $b,c$ nadal należące do tych samych przekrojów i nadal spełniające $q<a(b+c)$. Wybieramy dodatnie $\eta$ tak małe, że $q<a(b+c)-2\eta$. Wtedy $ab-\eta\in A\odot B$, $ac-\eta\in A\odot C$, a ich suma jest większa od $q$. Domknięcie w dół daje $q\in (A\odot B)+(A\odot C)$.
Odwrotnie, niech $q=u+v\ge0$, $u\in A\odot B$, $v\in A\odot C$. Istnieją dodatnie $a_1,a_2\in A$, $b\in B,c\in C$ z $u<a_1b$, $v<a_2c$; dla ujemnego $u$ lub $v$ wystarczy dowolna para dodatnich świadków. Kładziemy $a=\max(a_1,a_2)$, wtedy $q<a(b+c)$. To dowodzi drugiej inkluzji dla $q\ge0$; liczby ujemne należą do obu nieujemnych przekrojów. Gdy $A=0_D$, obie strony są zerem. Gdy $B=0_D$ lub $C=0_D$, obie strony redukują się odpowiednio do $A\odot C$ lub $A\odot B$.
\end{proof}
\begin{twierdzenie}[Uporządkowany pierścień]
\label{ma1:thm:cut-ring}
Istnieje dokładnie jedno mnożenie na $\R_D$ rozszerzające $\odot$ i czyniące $\R_D$ pierścieniem przemiennym z jedynką $1_D$. Porządek jest zgodny z dodawaniem i mnożeniem przez liczby dodatnie.
\end{twierdzenie}
\begin{proof}
Każdy element grupy $\R_D$ jest różnicą nieujemnych elementów, ponieważ $A\ge0_D$ daje $A=A-0_D$, a $A<0_D$ daje $A=0_D-(-A)$. Dla $A=P-Q$, $B=U-V$, gdzie $P,Q,U,V\ge0_D$, definiujemy
\[
 AB=P\odot U+Q\odot V-P\odot V-Q\odot U.
\]
Aby wykazać niezależność od rozkładu $A$, załóżmy $P-Q=P'-Q'$, czyli $P+Q'=P'+Q$. Dla dowolnego $W\ge0_D$ rozdzielność daje $P\odot W+Q'\odot W=P'\odot W+Q\odot W$, zatem $P\odot W-Q\odot W=P'\odot W-Q'\odot W$. Stosujemy to do $W=U,V$; niezależność od rozkładu $B$ wynika przez przemienność wzoru.
Rozdzielność wynika z dodawania rozkładów: $(P-Q)+(P'-Q')=(P+P')-(Q+Q')$ oraz rozdzielności $\odot$. Dla $A,B\ge0_D$ wybór rozkładów $A=A-0_D$, $B=B-0_D$ pokazuje, że $AB=A\odot B$. Dla dowolnego $A=P-Q$ wzór daje $A1_D=P-Q=A$. Przemienność wynika z przemienności $\odot$. Dla łączności zauważmy, że z rozdzielności oba wyrażenia $(AB)C$ i $A(BC)$ są addytywne osobno w każdym argumencie; po zapisaniu argumentów jako różnic nieujemnych elementów wystarczy porównać je na trójkach nieujemnych, gdzie zachodzi łączność $\odot$. Jednoznaczność wynika z konieczności rozwinięcia iloczynu różnic przez rozdzielność.
Jeśli $A<B$ i $C>0_D$, to $B-A>0_D$, a $(B-A)C>0_D$ przez poprzedni lemat. Stąd $AC<BC$. 
\end{proof}
\begin{twierdzenie}[Odwrotność i zanurzenie liczb wymiernych]
\label{ma1:thm:cut-field}
Pierścień $\R_D$ jest ciałem uporządkowanym, a odwzorowanie $r\mapsto D_r$ jest zanurzeniem ciał uporządkowanych.
\end{twierdzenie}
\begin{proof}
Dla $A>0_D$ definiujemy
\[
 B=\{q\in\Q:q<0\ \text{lub istnieje }u\in\Q\setminus A,\ u>0,\ q<1/u\}.
\]
Zbiór $B$ jest niepusty, domknięty w dół i bez maksimum. Wybierzmy dodatnie $a\in A$. Każde $u\notin A$ spełnia $u>a$, więc $1/a\notin B$. Zatem $B$ jest przekrojem dodatnim (zawiera np. $1/(2u)$ dla dowolnego dodatniego $u\notin A$).
Jeżeli $a\in A,b\in B$ są dodatnie, to $b<1/u$ dla pewnego $u\notin A$ i $a<u$, więc $ab<1$. Stąd $AB\subseteq1_D$. Niech $0\le q<1$. Wybierzmy $a_0>0$ w $A$ i $\eta>0$ wymierne z $\eta<(1-q)a_0$. Lemat \ref{ma1:lem:cut-bracket} daje $a\in A$, $u\notin A$, $u-a<\eta$. Możemy zastąpić $a$ przez $\max(a,a_0)$; nadal $u-a<\eta$ i $u>a\ge a_0$. Wówczas $1-a/u=(u-a)/u<1-q$, czyli $q<a/u$. Wybieramy dodatnie wymierne $b$ z $q/a<b<1/u$. Mamy $b\in B$, $q<ab$, więc $q\in AB$. Liczby ujemne również należą do $AB$, zatem $AB=1_D$. Dla $A<0_D$ mamy $-A>0_D$. Jeśli $C$ jest odwrotnością $-A$, to z rozdzielności w pierścieniu $A(-C)=(-A)C=1_D$. Ponieważ $0_D\ne1_D$, otrzymaliśmy ciało.
Zgodność zanurzenia z dodawaniem wykazano wcześniej. Dla dodatnich $r,s\in\Q$ z $a<r$, $b<s$, $a,b>0$, wynika $ab<rs$. Jeżeli $0\le q<rs$, wybieramy dodatnie $a$ z $q/s<a<r$, a następnie dodatnie $b$ z $q/a<b<s$. Wówczas $q<ab$. Stąd $D_rD_s=D_{rs}$. Dla argumentu zerowego wniosek wynika z $D_0=0_D$ i prawa zera. Ponieważ $D_r+D_{-r}=D_0$, mamy $D_{-r}=-D_r$. Dla argumentów ujemnych stosujemy zatem prawa znaków w pierścieniu i już wykazaną równość dla dodatnich argumentów. Zanurzenie zachowuje jedynkę i porządek, jest iniektywne na mocy twierdzenia \ref{ma1:thm:cut-complete}.
\end{proof}
Od tej chwili utożsamiamy $\Q$ z obrazem w $\R_D$. Własność Archimedesa zachodzi także w $\R_D$: dla $A\in\R_D$ wybieramy $u\notin A$ i liczbę naturalną $n>u$; wtedy $A<D_n$. Dla każdego $x>0$ w $\R_D$ istnieje dodatnia liczba wymierna mniejsza od $x$ dzięki gęstości z twierdzenia \ref{ma1:thm:cut-complete}.

\section{Ciągi Cauchy'ego liczb wymiernych}
\label{ma1:sec:rational-cauchy}

Przekrój opisuje liczbę przez jej położenie względem wszystkich liczb wymiernych. Możemy jednak opisać ją także przez coraz dokładniejsze przybliżenia wymierne. Aby przybliżenia określały jedną liczbę, ich wzajemne odległości powinny być dowolnie małe, gdy oba indeksy są dostatecznie duże. Zanim sformalizujemy tę intuicję, potrzebujemy kilku własności wartości bezwzględnej.

Niech $F$ będzie ciałem uporządkowanym (w tej sekcji wystarczy $F=\Q$). Dla $x\in F$ przyjmujemy $|x|=\max(x,-x)$. Wówczas $|x|\ge0$, przy czym $|x|=0$ wtedy i tylko wtedy, gdy $x=0$. Dla $x,y\in F$ mamy
\[
 -(|x|+|y|)\le x+y\le |x|+|y|,
\]
gdyż $-|x|\le x\le |x|$ oraz $-|y|\le y\le |y|$. Zatem $|x+y|\le|x|+|y|$ (nierówność trójkąta).
Jeśli $x,y$ mają ten sam znak, to $xy\ge0$ i $|x||y|=xy$; jeśli mają przeciwne znaki, to $xy\le0$ i $|x||y|=-xy$. Oba wzory są zgodne dla argumentu zerowego, a więc $|xy|=|x||y|$. W szczególności dla $y\ne0$ mamy $|1/y|=1/|y|$, ponieważ $|y||1/y|=1$.
Z nierówności trójkąta otrzymujemy $|x|\le|x-y|+|y|$ oraz $|y|\le|x-y|+|x|$, skąd
\[
 \bigl||x|-|y|\bigr|\le|x-y|.
\]
\begin{definicja}
\label{ma1:def:rational-cauchy}
Niech $(a_n)_{n\ge1}$ będzie ciągiem liczb wymiernych, czyli odwzorowaniem $\N_+\to\Q$. Nazywamy go \emph{ciągiem Cauchy'ego}, jeśli dla każdego wymiernego $\varepsilon>0$ istnieje $N\in\N_+$ takie, że dla wszystkich $m,n\ge N$ zachodzi $|a_n-a_m|<\varepsilon$. Ciąg $(a_n)_{n\ge1}$ liczb wymiernych nazywamy zbieżnym do zera, jeśli dla każdego wymiernego $\varepsilon>0$ istnieje $N\in\N_+$ takie, że $|a_n|<\varepsilon$ dla wszystkich $n\ge N$. W tym rozdziale równoważnie mówimy o ciągu zerowym w granicy.
\end{definicja}
\begin{przyklad}[Ciąg Cauchy'ego i ciąg ograniczony niespełniający warunku Cauchy'ego]
Dla $a_n=1/n$ i wymiernego $\varepsilon>0$ wybieramy $N\in\N_+$ z $N>2/\varepsilon$. Dla $n,m\ge N$ mamy $|a_n-a_m|\le1/n+1/m\le2/N<\varepsilon$, więc $(a_n)$ jest ciągiem Cauchy'ego. Ponadto $|a_n|\le1/N<\varepsilon$, czyli jest zbieżny do zera. Ciąg $b_n=(-1)^n$ nie jest ciągiem Cauchy'ego: dla $\varepsilon=1$ i dowolnego $N\in\N_+$ indeksy $n=N$ oraz $m=N+1$ spełniają $|b_n-b_m|=2$. Sama ograniczoność nie wystarcza zatem do spełnienia warunku Cauchy'ego.
\end{przyklad}

\begin{lemat}
\label{ma1:lem:cauchy-algebra}
Każdy wymierny ciąg Cauchy'ego jest ograniczony. Suma, różnica i iloczyn dwóch takich ciągów są ciągami Cauchy'ego. Suma ciągów zerowych w granicy oraz iloczyn ciągu zerowego w granicy z ograniczonym ciągiem są zerowe w granicy.
\end{lemat}
\begin{proof}
Dla $\varepsilon=1$ wybieramy $N$ z definicji, wtedy $|a_n|\le|a_N|+1$ dla $n\ge N$; maksimum tej liczby i modułów początkowych wyrazów ogranicza cały ciąg. Niech $\varepsilon>0$ będzie wymierne. Dla sumy i różnicy wybieramy wspólny indeks $N$ z warunku Cauchy'ego obu ciągów z dokładnością $\varepsilon/2$ (maksimum dwóch uzyskanych indeksów). Dla $m,n\ge N$ mamy
\[
 |(a_n\pm b_n)-(a_m\pm b_m)|
 \le|a_n-a_m|+|b_n-b_m|<\varepsilon.
\]
Dla iloczynu wybieramy wspólne wymierne ograniczenie $M\ge1$ modułów obu ciągów. Zauważmy, że
\[
\begin{aligned}
 |a_nb_n-a_mb_m|
 &=|a_n(b_n-b_m)+(a_n-a_m)b_m|\\
 &\le |a_n||b_n-b_m|+|b_m||a_n-a_m|\\
 &\le M(|a_n-a_m|+|b_n-b_m|).
\end{aligned}
\]
Wybierając wspólny indeks z warunku Cauchy'ego z dokładnością $\varepsilon/(2M)$, otrzymujemy dla wszystkich późniejszych $m,n$ oszacowanie $|a_nb_n-a_mb_m|<\varepsilon$.

Jeżeli oba ciągi są zbieżne do zera, to od pewnego miejsca $|a_n|,|b_n|<\varepsilon/2$, a stąd $|a_n+b_n|<\varepsilon$. Jeżeli $a_n$ jest zbieżny do zera, a $|b_n|\le M$ dla wymiernego $M\ge1$, to od pewnego miejsca $|a_n|<\varepsilon/M$ i $|a_nb_n|\le M|a_n|<\varepsilon$.
\end{proof}
\begin{definicja}
Na zbiorze wymiernych ciągów Cauchy'ego definiujemy $a\sim b$, gdy $(a_n-b_n)$ jest zerowy w granicy. Oznaczamy zbiór klas przez $\R_C$ i klasę ciągu przez $[a]$.
\end{definicja}
\begin{twierdzenie}
\label{ma1:thm:cauchy-ring}
Relacja $\sim$ jest relacją równoważności. Działania $[a]+[b]=[(a_n+b_n)]$, $[a][b]=[(a_nb_n)]$ czynią $\R_C$ pierścieniem przemiennym z jedynką, z zanurzeniem $\Q$ przez ciągi stałe.
\end{twierdzenie}
\begin{proof}
Zwrotność i symetria wynikają z modułu, przechodniość z nierówności trójkąta. Operacje zachowują warunek Cauchy'ego przez lemat \ref{ma1:lem:cauchy-algebra}. Jeśli $a\sim a'$ i $b\sim b'$, to różnica sum jest sumą dwóch ciągów zerowych w granicy, a
$a_nb_n-a'_nb'_n=(a_n-a'_n)b_n+a'_n(b_n-b'_n)$
jest również zerowa w granicy. Prawa pierścienia dziedziczymy wyraz po wyrazie; zero i jedynka to klasy ciągów stałych $0$ i $1$, różne, bo ciąg stały $1$ nie jest zerowy w granicy. Negację otrzymujemy wyraz po wyrazie. Różnica dwóch różnych ciągów stałych jest stałą niezerową i nie jest zerowa w granicy, co daje iniektywność zanurzenia.
\end{proof}

\begin{twierdzenie}[Odwrotność klasy ciągu Cauchy'ego]
\label{ma1:thm:cauchy-inverse}
Niech $a=(a_n)_{n\ge1}$ będzie wymiernym ciągiem Cauchy'ego o niezerowej klasie $[a]\in\R_C$. Istnieją $c\in\Q$, $c>0$, oraz $N\in\N_+$ takie, że $|a_n|\ge c$ dla każdego $n\ge N$. Ciąg $b$ określony przez $b_n=1/a_n$ dla $n\ge N$ i $b_n=0$ dla $n<N$ jest ciągiem Cauchy'ego oraz $[a][b]=1$. W szczególności $\R_C$ jest ciałem niezależnie od konstrukcji Dedekinda.
\end{twierdzenie}
\begin{proof}
\emph{Ograniczenie modułu od dołu.} Ponieważ $[a]\ne0$, ciąg $a$ nie jest zbieżny do zera. Istnieje więc wymierne $\eta>0$ takie, że dla każdego $K\in\N_+$ istnieje $k\ge K$ z $|a_k|\ge\eta$. Z warunku Cauchy'ego wybieramy $N$ tak, aby $|a_n-a_m|<\eta/2$ dla wszystkich $n,m\ge N$. Wybieramy indeks $k\ge N$ z $|a_k|\ge\eta$. Dla każdego $n\ge N$ mamy
\[
 |a_n|\ge |a_k|-|a_n-a_k|>\eta/2.
\]
Kładziemy $c=\eta/2$. Wszystkie wyrazy $a_n$ dla $n\ge N$ są więc niezerowe.

\emph{Warunek Cauchy'ego dla odwrotności.} Niech $\varepsilon>0$ będzie wymierne. Wybieramy $K\ge N$ tak, aby $|a_n-a_m|<c^2\varepsilon$ dla $n,m\ge K$. Wtedy
\[
 |b_n-b_m|=\frac{|a_m-a_n|}{|a_na_m|}
 \le\frac{|a_m-a_n|}{c^2}<\varepsilon.
\]
Zatem $b$ jest ciągiem Cauchy'ego.

\emph{Iloczyn klas i jednoznaczność.} Dla $n\ge N$ zachodzi $a_nb_n=1$, więc różnica ciągu $(a_nb_n)$ i ciągu stałego $1$ jest zbieżna do zera: dla każdego wymiernego $\varepsilon>0$ jej moduł wynosi $0<\varepsilon$ od indeksu $N$. Na mocy definicji działań na klasach $[a][b]=1$. Odwrotność klasy jest jednoznaczna: jeśli $xy=xz=1$, to $y=y(xz)=(yx)z=z$. W konsekwencji wybór reprezentanta klasy i wartości skończenie wielu początkowych wyrazów nie zmienia klasy odwrotnej.
\end{proof}

\section{Równoważność konstrukcji}
\label{ma1:sec:real-equivalence}
\begin{lemat}[Zbieżność wymiernych ciągów Cauchy'ego w $\R_D$]
\label{ma1:lem:cauchy-limit}
Dla każdego wymiernego ciągu Cauchy'ego $(a_n)_{n\ge1}$ istnieje dokładnie jedno $L\in\R_D$ takie, że dla każdego wymiernego $\varepsilon>0$ istnieje $N\in\N_+$ spełniające $|D_{a_n}-L|<D_\varepsilon$ dla wszystkich $n\ge N$. Po utożsamieniu $\Q$ z jej obrazem jest to warunek $|a_n-L|<\varepsilon$.
\end{lemat}
\begin{proof}
Utożsamiamy $a_n$ z przekrojem $D_{a_n}$. Ograniczoność ciągu zapewnia istnienie $s_N=\inf\{a_n:n\ge N\}$ i $t_N=\sup\{a_n:n\ge N\}$ w $\R_D$; kres dolny istnieje, gdyż dla niepustego zbioru $E$ ograniczonego z dołu zbiór $-E=\{-x:x\in E\}$ jest ograniczony z góry, a $-\sup(-E)$ jest ograniczeniem dolnym $E$ większym lub równym każdemu innemu ograniczeniu dolnemu (dodawanie elementów przeciwnych odwraca porządek). Ciąg $s_N$ jest niemalejący i ograniczony z góry, więc $L=\sup\{s_N:N\ge1\}$ istnieje. Ustalmy $N$. Dla każdego $K\ge1$ wybieramy $n\ge\max(N,K)$; mamy $s_K\le a_n\le t_N$. Liczba $t_N$ jest zatem ograniczeniem górnym wszystkich $s_K$, więc $L\le t_N$ z definicji kresu górnego. Zatem $s_N\le L\le t_N$.
Dla wymiernego $\varepsilon>0$ wybierzmy $N$ z warunku Cauchy'ego dla $\varepsilon/2$. Dla $n,m\ge N$ mamy $a_n-\varepsilon/2\le a_m\le a_n+\varepsilon/2$, co po wzięciu kresów daje $s_N\ge a_n-\varepsilon/2$ i $t_N\le a_n+\varepsilon/2$. W konsekwencji $|a_n-L|\le\varepsilon/2<\varepsilon$.
Jeśli $L'$ również spełnia żądany warunek, a $L\ne L'$, wybieramy wymierne $0<\varepsilon<|L-L'|/2$ dzięki gęstości $\Q$ w $\R_D$. Dla dostatecznie dużego $n$ nierówność trójkąta daje $|L-L'|<2\varepsilon$, sprzeczność.
\end{proof}
\begin{twierdzenie}[Izomorfizm konstrukcji]
\label{ma1:thm:real-isomorphism}
Odwzorowanie $\Phi\colon\R_C\to\R_D$ przypisujące klasie ciągu jego granicę z lematu \ref{ma1:lem:cauchy-limit} jest izomorfizmem pierścieni. Jest więc także izomorfizmem ciał (por. twierdzenie \ref{ma1:thm:cauchy-inverse}). Porządek przeniesiony z $\R_D$ charakteryzuje warunek
\[
 [a]<[b]\quad\Longleftrightarrow\quad
 \exists\varepsilon\in\Q,\ \varepsilon>0,\ \exists N\ge1\
 \forall n\ge N:\ b_n-a_n>\varepsilon .
\]
Obie konstrukcje dają ciało uporządkowane zupełne i zgodne zanurzenia $\Q$.
\end{twierdzenie}
\begin{proof}
Jeśli $a\sim b$, nierówność trójkąta pokazuje, że granica $a$ jest także granicą $b$: dla żądanej dokładności $\varepsilon$ stosujemy oba warunki z $\varepsilon/2$. Odwrotnie, jeśli granice są równe, dla dużych $n$ mamy $|a_n-b_n|\le|a_n-L|+|b_n-L|<\varepsilon$. To daje dobre określenie i iniektywność.
Niech $L=\Phi([a])$, $K=\Phi([b])$ i niech $\varepsilon>0$ będzie wymierne. Dla dostatecznie dużych $n$ oba błędy $|a_n-L|$ i $|b_n-K|$ są mniejsze od $\varepsilon/2$, a zatem
$|(a_n+b_n)-(L+K)|\le|a_n-L|+|b_n-K|<\varepsilon$.
Jednoznaczność granicy daje $\Phi([a]+[b])=L+K$.

Dla iloczynu wybieramy wymierne $M\ge1$ ograniczające wszystkie $|a_n|$ oraz $|K|$ (dla $K$ używamy własności Archimedesa w $\R_D$). Ponieważ
\[
\begin{aligned}
 |a_nb_n-LK|
 &=|a_n(b_n-K)+(a_n-L)K|\\
 &\le |a_n||b_n-K|+|K||a_n-L|\\
 &\le M(|b_n-K|+|a_n-L|),
\end{aligned}
\]
dla dostatecznie dużych $n$, wybranych przy dokładności $\varepsilon/(2M)$, ostatnia wielkość jest mniejsza od $\varepsilon$. Jednoznaczność granicy daje $\Phi([a][b])=LK$. Dla każdego $r\in\Q$ ciąg stały o wartości $r$ ma granicę $D_r$, ponieważ jego różnica z $D_r$ jest zerem w $\R_D$. W szczególności jedynka i zero są zachowane, a oba zanurzenia $\Q$ są zgodne.
Dla dowolnego $L\in\R_D$ wybieramy dla każdego $n\ge1$ liczbę wymierną $a_n$ z $L-1/n<a_n<L+1/n$ (gęstość). Wówczas $|a_n-L|<1/n$, a własność Archimedesa daje dla każdego wymiernego $\varepsilon>0$ istnienie $N$ z $2/N<\varepsilon$. Stąd dla $n,m\ge N$ zachodzi $|a_n-a_m|\le |a_n-L|+|a_m-L|<1/n+1/m\le2/N<\varepsilon$. Ponadto $|a_n-L|<1/n\le1/N<\varepsilon$, więc granicą jest $L$, zatem $\Phi$ jest surjekcją. Możemy dokonać wyboru $a_n$ deterministycznie, biorąc pierwszy odpowiedni element ustalonego wyliczenia $\Q$ (pary licznika i dodatniego mianownika można uporządkować rosnąco według sumy modułu licznika i mianownika, a następnie rozstrzygać remisy według licznika).
Jeżeli $K-L>0$, wybieramy wymierne $\varepsilon>0$ z $3\varepsilon<K-L$; dla dużych $n$ mamy $|b_n-K|<\varepsilon$, $|a_n-L|<\varepsilon$, więc $b_n-a_n>K-L-2\varepsilon>\varepsilon$. Jeśli odwrotnie zachodzi warunek prawej strony, a $K-L\le0$, to dla dużych $n$ oba błędy są mniejsze niż $\varepsilon/2$, co daje $b_n-a_n<\varepsilon$, sprzeczność. Zatem porządek opisany w tezie jest dobrze określony i zgodny z $\Phi$. Aby przenieść zupełność, niech $\mathcal B\subseteq\R_C$ będzie niepusta i ograniczona z góry przez $U$. Obraz $\Phi(\mathcal B)$ jest niepusty i ograniczony z góry przez $\Phi(U)$. Istnieje więc $S=\sup\Phi(\mathcal B)$ w $\R_D$. Z surjektywności $\Phi$ istnieje $T=\Phi^{-1}(S)$. Ponieważ $\Phi$ zachowuje i odzwierciedla porządek, dla każdego $B\in\mathcal B$ mamy $B\le T$. Jeżeli $V$ jest dowolnym ograniczeniem górnym $\mathcal B$, to $\Phi(V)$ jest ograniczeniem górnym obrazu, więc $S\le\Phi(V)$, czyli $T\le V$. W konsekwencji $T=\sup\mathcal B$.
\end{proof}
\begin{przyklad}
Przekrój $A=\{q\in\Q:q<0\ \text{lub }q^2<2\}$ jest poprawny: dla $q\ge0$ z $q^2<2$ liczba $(2q+2)/(q+2)$ jest większa od $q$ i nadal ma kwadrat mniejszy od $2$, jak w rozdziale \ref{ma1:ch:number-systems}. Dla $q<0$ można wybrać większą liczbę ujemną. Domknięcie w dół wynika z monotoniczności kwadratu na liczbach nieujemnych. Zawiera $1$, nie zawiera $2$. Kres górny zbioru $\{D_q:q\in A\}$ w $\R_D$ jest równy $A$: zgodnie z twierdzeniem \ref{ma1:thm:cut-complete} jest to $\bigcup_{q\in A}D_q$. Inkluzja tej sumy w $A$ wynika z domknięcia w dół, a odwrotna z braku elementu największego w $A$. W następnym rozdziale, po ogólnym twierdzeniu o pierwiastkach dodatnich, zidentyfikujemy go z $\sqrt2$.
\end{przyklad}
\begin{uwaga}
Odtąd piszemy $\R$ dla skonstruowanego ciała. Wykazaliśmy zupełność porządkową, nie zakładając wcześniej ogólnej teorii przestrzeni metrycznych. W dalszych rozdziałach rozwiniemy tę teorię, już z odległościami o wartościach w $\R$.
\end{uwaga}
